\section{Proof of the extended KRR--S theorem}\label{app:krrs-analytic-extension}

We first record the results from~\cite{kenyon2014entropy} needed for the
proof of \Cref{thm:krrs-analytic-extension}, restricting throughout to
$d$-regular graphs.
Their Theorem~1.1 describes the optimizer parameters only for
\(\vartheta:=\tau-\eps^m>0\).  The rest of this section constructs their
two-sided analytic continuation and verifies that it agrees with those
optimizer parameters where they are defined.

\begin{theorem}[{KRR--S bipodality for regular graphs, \cite[Theorem~1.1]{kenyon2014entropy}}]
\label{thm:krrs-bipodality}
Let \(H\) be a \(d\)-regular graph with \(d\ge2\) and \(m\ge2\) edges, and let
\(\eps\in(0,1)\setminus\{(d-1)/d\}\).  There exists
\(\tau_0(\eps)>\eps^m\) such that, whenever
\(\eps^m<\tau<\tau_0(\eps)\), the entropy maximizer subject to
\[
  e(W)=\eps,
  \qquad
  t(H,W)=\tau
\]
is unique up to relabeling and is bipodal.  Let $c$ be the size of the smaller block,
and let $q_{11},q_{12},q_{22}$ denote the edge probabilities within the smaller block (of size $c$),
between the two blocks, and within the larger block (of size $1-c$), respectively.
Then, its parameters \((c, q_{11},q_{12},q_{22})\) are jointly
real-analytic in \((\eps,\tau)\) on the region
\[
  \left\{(\eps,\tau):
    \eps\ne\frac{d-1}{d},\quad
    \eps^m<\tau<\tau_0(\eps)\right\}.
\]
For each fixed nonexceptional \(\eps\), as \(\tau\downarrow\eps^m\),
\[
  q_{11}\longrightarrow a_d(\eps),
  \qquad
  q_{12}\longrightarrow\zeta_d(\eps),
  \qquad
  q_{22}\longrightarrow\eps,
  \qquad
  c=O_{H,\eps}(\tau-\eps^m),
\]
where \(a_d(\eps)\in(0,1)\) and \(\zeta_d(\eps)\ne\eps\).
\end{theorem}

\begin{theorem}[{KRR--S cross density for regular graphs, \cite[Theorem~3.3]{kenyon2014entropy}}]
\label{thm:krrs-cross-density}
For every \(d\ge2\) and \(\eps\in(0,1)\), the function
\(z\mapsto\psi_d(\eps,z)\) has a unique critical point, denoted
\(\zeta_d(\eps)\), at which it attains its maximum on \((0,1)\).
The map \(\zeta_d:(0,1)\to(0,1)\) is strictly decreasing, and is an involution, i.e.,
$\eps=\zeta_d(\zeta_d(\eps))$.  Its unique fixed point is \((d-1)/d\).
\end{theorem}

\begin{remark}\label{rmk:krrs-omitted-identity}
We omit the identity \(S_0'(q_{11})=2S_0'(q_{12})-S_0'(q_{22})\)
asserted in \cite[Theorem~1.1]{kenyon2014entropy} and the claim in
\cite[Theorem~3.3]{kenyon2014entropy} that $\zeta_d'$ is nowhere zero.
Neither assertion is used in our subsequent arguments.
In place of the identity for $q_{11}$, we use only that its limiting value
$a_d(\eps)$ lies in $(0,1)$.
\end{remark}

Before proving \Cref{thm:krrs-analytic-extension}, we clarify what must be added to
\Cref{thm:krrs-bipodality}.  Although the latter states the threshold
\(\tau_0(\eps)\) pointwise, its joint-analyticity assertion implies that, for
every nonexceptional \(\eps_0\), there are a neighborhood \(U\) of \(\eps_0\)
and \(\Delta>0\) such that the theorem statement holds uniformly for
\(U\times(0,\Delta)\).  It does not, however, provide an analytic extension
through \(\vartheta=0\).  The main purpose of the following proof is to construct this
two-sided jointly analytic extension, which is needed for the boundary
analysis in
\Cref{thm:positive-second-variation,thm:nonexceptional-endpoint}.

\begin{proof}[Proof of \Cref{thm:krrs-analytic-extension}]
Let $\eps_*:=r_*=(d-1)/d$
and fix
\[
  \eps_0\in(0,1)\setminus\{\eps_*\},
  \qquad
  a_0:=a_d(\eps_0).
  \qquad
  z_0:=\zeta_d(\eps_0),
\]
By \Cref{thm:krrs-bipodality}, \(z_0\ne\eps_0\) and the parameters converge to
\begin{equation}\label{eq:krrs-parameter-limits}
  (q_{11},q_{12},q_{22},c)\longrightarrow(a_0,z_0,\eps_0,0).
\end{equation}
The local Taylor remainders below are uniform on sufficiently small fixed
neighborhoods of this base point, chosen in terms of $H$ and $\eps_0$.

\paragraph{1. Nondegeneracy of the scalar maximizer.}
For fixed \(\eps\), put
\begin{align}
  \mathcal N_\eps(z)
  &:=2\bigl[S_0(z)-S_0(\eps)-S_0'(\eps)(z-\eps)\bigr],\label{eq:krrs-entropy-remainder}\\
  \mathcal D_\eps(z)
  &:=z^d-\eps^d-d\eps^{d-1}(z-\eps).\label{eq:krrs-moment-remainder}
\end{align}
Thus, \(\psi_d(\eps,z)=\mathcal N_\eps(z)/\mathcal D_\eps(z)\) when
\(z\ne\eps\).  We claim that
\begin{equation}\label{eq:krrs-scalar-nondegeneracy}
  \partial_z^2\psi_d(\eps_0,z_0)<0.
\end{equation}
To see this, \(z_0\) is an interior maximizer, so its second derivative is
nonpositive.  Suppose for contradiction that it vanishes.  The third
derivative must then vanish as well, since an interior local maximum cannot
have a first nonzero derivative of odd order.  At \((\eps_0,z_0)\) we have
\(z_0\ne\eps_0\), and hence the identity
\(\mathcal N_{\eps_0}=\psi_d(\eps_0,\cdot)\mathcal D_{\eps_0}\) holds.  Since the first three derivatives of \(\psi_d\)
at \(z_0\) would vanish, by differentiating this identity we would obtain
\begin{equation}\label{eq:krrs-ratio-derivatives}
  \mathcal N_{\eps_0}''(z_0)
    =\psi_d(\eps_0,z_0)\mathcal D_{\eps_0}''(z_0),
  \qquad
  \mathcal N_{\eps_0}'''(z_0)
    =\psi_d(\eps_0,z_0)\mathcal D_{\eps_0}'''(z_0).
\end{equation}
The relevant derivatives are
\begin{align*}
  \mathcal N_\eps''(z)&=-\frac1{z(1-z)},
  &
  \mathcal N_\eps'''(z)&=\frac{1-2z}{z^2(1-z)^2},\\
  \mathcal D_\eps''(z)&=d(d-1)z^{d-2},
  &
  \mathcal D_\eps'''(z)&=d(d-1)(d-2)z^{d-3}.
\end{align*}
If \(d=2\), then \(\mathcal D_{\eps_0}'''\equiv0\), so the second identity
in \eqref{eq:krrs-ratio-derivatives} gives
\(\mathcal N_{\eps_0}'''(z_0)=0\), hence \(z_0=1/2=\eps_*\).  If
\(d\ge3\), divide the second identity in
\eqref{eq:krrs-ratio-derivatives} by the first; the division is legitimate
because both sides of the first equal
\(\mathcal N_{\eps_0}''(z_0)=-1/(z_0(1-z_0))\ne0\), and it yields
\[
  \frac{2z_0-1}{z_0(1-z_0)}=\frac{d-2}{z_0},
\]
hence again \(z_0=(d-1)/d=\eps_*\).  In either case, this implies $\eps_0 = (d-1)/d$,
contrary to the choice of \(\eps_0\).  This proves
\eqref{eq:krrs-scalar-nondegeneracy}.

\paragraph{2. An analytic parametrization of the two constraints.}
Write a bipodal graphon, with the first pode of size \(c\), as
\[
  W(a,b,q,c),
  \qquad
  (a,b,q)=(q_{11},q_{12},q_{22}).
\]
Its edge density is
\[
  \mathcal E(a,b,q,c)
  =c^2a+2c(1-c)b+(1-c)^2q.
\]
For prescribed edge density \(\eps\), solve this equation exactly for \(q\):
\begin{equation}\label{eq:krrs-edge-constraint}
  q=Q(\eps,a,b,c)
  :=\frac{\eps-c^2a-2c(1-c)b}{(1-c)^2}.
\end{equation}
The function \(Q\) is real-analytic near
\((\eps_0,a_0,z_0,0)\), and
\begin{equation}\label{eq:krrs-large-block-expansion}
  Q(\eps,a,b,c)
  =\eps+2(\eps-b)c+(3\eps-2b-a)c^2+O(c^3).
\end{equation}
Here the remainder has an absolute constant, uniformly for
$(\eps,a,b)\in[0,1]^3$ and $|c|\le1/2$.
Define the edge-constrained \(H\)-density and entropy by
\begin{align*}
  \widehat{\mathcal T}(\eps,a,b,c)
    &:=t\bigl(H,W(a,b,Q(\eps,a,b,c),c)\bigr),\\
  \widehat{\mathcal S}(\eps,a,b,c)
    &:=c^2S_0(a)+2c(1-c)S_0(b)
       +(1-c)^2S_0\bigl(Q(\eps,a,b,c)\bigr).
\end{align*}

Let \(n_1\) and \(n_d\) be the numbers of vertices of \(H\) of degrees
\(1\) and \(d\), respectively, and put \(v=n_1+n_d\).  Then
\(n_1+dn_d=2m\), and \(n_d>0\).  Expanding the homomorphism-density sum
according to the set of vertices assigned to the first pode gives
\begin{align}
  \left.\partial_c\widehat{\mathcal T}(\eps,a,b,c)\right|_{c=0}
  &=\sum_{x\in V(H)}b^{\deg(x)}\eps^{m-\deg(x)}
    -v\eps^m+2m\eps^{m-1}(\eps-b)\notag\\
  &=n_d\eps^{m-d}
      \bigl[b^d-\eps^d-d\eps^{d-1}(b-\eps)\bigr].
      \label{eq:krrs-density-derivative}
\end{align}
To verify the first equality, consider first an assignment in which \(x\) is
the unique vertex in the first pode.  Its contribution to
\(\widehat{\mathcal T}\) is
\[
  c(1-c)^{v-1}b^{\deg(x)}
  Q(\eps,a,b,c)^{m-\deg(x)},
\]
whose derivative at \(c=0\) is
\(b^{\deg(x)}\eps^{m-\deg(x)}\).  Summing over \(x\in V(H)\) gives the first
term.  If at least two vertices are assigned to the first pode, the
corresponding summand contains a factor \(c^k\) for some \(k\ge2\), while
all its other factors are analytic near \(c=0\); hence its derivative at
\(c=0\) vanishes.  Finally, the all-second-pode assignment contributes
\((1-c)^vQ(\eps,a,b,c)^m\).  Since \(Q(\eps,a,b,0)=\eps\) and
\(\partial_cQ|_{c=0}=2(\eps-b)\), its derivative at \(c=0\) is
\[
  -v\eps^m+m\eps^{m-1}\,2(\eps-b).
\]
This verifies the first equality in \eqref{eq:krrs-density-derivative}.
The second equality follows from the two identities for \(n_1,n_d\) above.

Set
\begin{equation}\label{eq:krrs-density-coefficient}
  \mathcal A(\eps,b)
  :=n_d\eps^{m-d}\mathcal D_\eps(b),
\end{equation}
which is equal to $\left.\partial_c\widehat{\mathcal T}(\eps,a,b,c)\right|_{c=0}$ by \eqref{eq:krrs-density-derivative}.
Strict convexity of \(z\mapsto z^d\) gives
\(\mathcal D_\eps(b)>0\) whenever \(b\ne\eps\).  In particular,
\(\mathcal A(\eps_0,z_0)>0\).  The analytic implicit function theorem,
applied to
\[
  \widehat{\mathcal T}(\eps,a,b,c)=\eps^m+\vartheta,
\]
gives neighborhoods of \((\eps_0,a_0,z_0,0)\) and a unique
real-analytic function
\begin{equation}\label{eq:krrs-density-constraint}
  c=C(\eps,a,b,\vartheta),
  \qquad
  C(\eps,a,b,0)=0,
\end{equation}
that solves the \(H\)-density constraint.  After shrinking the
neighborhood, \(Q\) and \(C\) take values in \((0,1)\) whenever
\(\vartheta>0\) and \((\eps,a,b)\) is sufficiently close to the base point:
for \(Q\), this follows by continuity and
\(Q(\eps_0,a_0,z_0,0)=\eps_0\in(0,1)\), while for \(C\) it follows from
\(C(\eps,a,b,0)=0\) together with
\(\partial_\vartheta C(\eps,a,b,0)=1/\mathcal A(\eps,b)>0\), obtained by
differentiating the constraint in \(\vartheta\).

Substitute \eqref{eq:krrs-density-constraint} into the entropy and write
\begin{equation}\label{eq:krrs-constrained-entropy}
  \Sigma(\eps,a,b,\vartheta)
  :=\widehat{\mathcal S}
       \bigl(\eps,a,b,C(\eps,a,b,\vartheta)\bigr).
\end{equation}
Since \(\Sigma(\eps,a,b,0)=S_0(\eps)\), there exists a real-analytic function \(L\) such that
\begin{equation}\label{eq:krrs-entropy-quotient}
  \Sigma(\eps,a,b,\vartheta)
  =S_0(\eps)+\vartheta L(\eps,a,b,\vartheta),
  \qquad
  L(\eps,a,b,\vartheta)
  :=\int_0^1
     \partial_\vartheta\Sigma(\eps,a,b,t\vartheta)\,dt.
\end{equation}
For \(\vartheta\ne0\), \(L\) agrees with
\((\Sigma-S_0(\eps))/\vartheta\), while the integral formula defines its
real-analytic extension to \(\vartheta=0\).

Using \(Q(\eps,a,b,0)=\eps\) and
\(\partial_cQ(\eps,a,b,0)=2(\eps-b)\), since the term \(c^2S_0(a)\) makes no contribution to the first derivative, direct differentiation gives
\begin{equation}\label{eq:krrs-entropy-derivative}
  \begin{aligned}
  \left.\partial_c\widehat{\mathcal S}(\eps,a,b,c)\right|_{c=0}
  &=2S_0(b)-2S_0(\eps)+2S_0'(\eps)(\eps-b)\\
  &=2\bigl[S_0(b)-S_0(\eps)-S_0'(\eps)(b-\eps)\bigr]\\
  &=\mathcal N_\eps(b).
  \end{aligned}
\end{equation}

Differentiating the \(H\)-density constraint with respect to \(\vartheta\) at
\(\vartheta=0\), and using
\eqref{eq:krrs-density-derivative} and \eqref{eq:krrs-density-coefficient}, gives
\[
  \mathcal A(\eps,b)\,\partial_\vartheta C(\eps,a,b,0)=1,
  \qquad
  \partial_\vartheta C(\eps,a,b,0)=\frac{1}{\mathcal A(\eps,b)}.
\]
Moreover, \eqref{eq:krrs-entropy-quotient} gives
\(L(\eps,a,b,0)=\partial_\vartheta\Sigma(\eps,a,b,0)\).  Hence, the chain rule
and \eqref{eq:krrs-entropy-derivative} yield the boundary identity
\begin{equation}\label{eq:krrs-quotient-boundary}
  \begin{aligned}
  L(\eps,a,b,0)
  &=\left.\partial_c\widehat{\mathcal S}(\eps,a,b,c)\right|_{c=0}
    \partial_\vartheta C(\eps,a,b,0)\\
  &=\frac{\mathcal N_\eps(b)}{\mathcal A(\eps,b)}
   =\frac{\psi_d(\eps,b)}{n_d\eps^{m-d}}.
  \end{aligned}
\end{equation}
Thus, the coefficient \(L(\eps,a,b,0)\) of \(\vartheta\) in the Taylor
expansion of the entropy \(\Sigma(\eps,a,b,\vartheta)\) about \(\vartheta=0\) is
independent of \(a=q_{11}\).

\paragraph{3. Construction of a two-sided analytic solution of the
stationarity equations.}
We next retain exactly the second-order information needed to recover the
otherwise invisible parameter \(a=q_{11}\).

We use the following elementary analytic factorization fact.  Let \(f\) be
real-analytic near a point of the hyperplane \(\{x_j=0\}\) and suppose
\(\partial_{x_j}^{\,i}f=0\) on that hyperplane for \(0\le i<k\).  Then
\(f=x_j^{k}g\) with \(g\) real-analytic.  For \(k=1\), one may take
\[
  g(x)=\int_0^1
    (\partial_jf)(x_1,\dots,tx_j,\dots,x_n)\,dt.
\]
The general case follows by iteration.  We apply this factorization in the
\(c\)-variable with \(k=3\).  For each of \(\widehat{\mathcal T}\) and
\(\widehat{\mathcal S}\), the difference from its second-order Taylor
polynomial at \(c=0\) has vanishing \(c\)-derivatives of orders \(0,1,2\) on
\(\{c=0\}\).  Hence, these differences can be written as
\(c^3R_{\mathcal T}\) and \(c^3R_{\mathcal S}\), respectively, with
\(R_{\mathcal T}\) and \(R_{\mathcal S}\) real-analytic.  Writing the
second-order Taylor coefficients in the form justified below gives
real-analytic functions \(B_0,B_1,R_0,R_1\), depending on \((\eps,b)\), such
that
\begin{align}
  \widehat{\mathcal T}(\eps,a,b,c)
  &=\eps^m+\mathcal A(\eps,b)c
    +\bigl[B_0(\eps,b)+B_1(\eps,b)a\bigr]c^2
    +c^3R_{\mathcal T}(\eps,a,b,c),
    \label{eq:krrs-density-expansion}\\
  \widehat{\mathcal S}(\eps,a,b,c)
  &=S_0(\eps)+\mathcal N_\eps(b)c
    +\bigl[S_0(a)+R_0(\eps,b)+R_1(\eps,b)a\bigr]c^2
    +c^3R_{\mathcal S}(\eps,a,b,c).
    \label{eq:krrs-entropy-expansion}
\end{align}
We now verify the stated form of the second-order coefficients in
\eqref{eq:krrs-density-expansion} and \eqref{eq:krrs-entropy-expansion}.  By
\eqref{eq:krrs-large-block-expansion}, write
\[
  Q(\eps,a,b,c)=\eps+q_1c+q_2c^2+O(c^3),
  \qquad q_1=2(\eps-b),\quad q_2=3\eps-2b-a.
\]
Thus, \(q_1\) is independent of \(a\), whereas \(q_2\) is affine in \(a\).
In the homomorphism-density expansion, the contribution at order \(c^2\)
from assignments with no vertex in the first pode depends linearly on
\(q_2\), while all its remaining terms depend only on \(q_1\).  An assignment
with one vertex in the first pode already carries a factor \(c\), so its
coefficient at order \(c^2\) depends only on \(q_1\).  An assignment with
exactly two vertices in the first pode contains at most one factor \(a\),
because \(H\) is simple and those two vertices support at most one edge.
Finally, assignments with at least three vertices in the first pode are
divisible by \(c^3\).  Hence, the coefficient of \(c^2\) in
\(\widehat{\mathcal T}\) is affine in \(a\), which proves
\eqref{eq:krrs-density-expansion}.

For \(\widehat{\mathcal S}\), the term \(c^2S_0(a)\) is the only possible
nonlinear dependence on \(a\).  To see this, \(2c(1-c)S_0(b)\) is independent of
\(a\), and the coefficient of \(c^2\) in the remaining term is
\[
  [c^2]\,(1-c)^2S_0(Q)
  =S_0(\eps)-2S_0'(\eps)q_1+S_0'(\eps)q_2
    +\frac12S_0''(\eps)q_1^2.
\]
Since \(q_1\) is independent of \(a\) and \(q_2\) is affine in \(a\), this
coefficient is affine in \(a\).  This proves
\eqref{eq:krrs-entropy-expansion}.

Under the constraint
\(\widehat{\mathcal T}(\eps,a,b,c)=\eps^m+\vartheta\), solving
\eqref{eq:krrs-density-expansion} for
\(c=C(\eps,a,b,\vartheta)\) as a power series in \(\vartheta\) gives
\begin{equation}\label{eq:krrs-block-size-expansion}
  \begin{aligned}
    C(\eps,a,b,\vartheta)
    &=\frac{\vartheta}{\mathcal A(\eps,b)}
      -\frac{B_0(\eps,b)+B_1(\eps,b)a}
             {\mathcal A(\eps,b)^3}\,\vartheta^2
      +O_{H,\eps_0}(\vartheta^3).
  \end{aligned}
\end{equation}
Substituting \eqref{eq:krrs-block-size-expansion} into
\eqref{eq:krrs-entropy-expansion} and recalling
\eqref{eq:krrs-entropy-quotient} gives
\begin{equation}\label{eq:krrs-quotient-expansion}
  L(\eps,a,b,\vartheta)
  =\frac{\mathcal N_\eps(b)}{\mathcal A(\eps,b)}
    +\vartheta M(\eps,a,b)+O_{H,\eps_0}(\vartheta^2),
\end{equation}
where
\begin{equation}\label{eq:krrs-second-coefficient}
  M(\eps,a,b)
  =\frac{S_0(a)+R_0(\eps,b)+R_1(\eps,b)a}
         {\mathcal A(\eps,b)^2}
   -\frac{\mathcal N_\eps(b)
          [B_0(\eps,b)+B_1(\eps,b)a]}
         {\mathcal A(\eps,b)^3}.
\end{equation}
Consequently,
\begin{equation}\label{eq:krrs-coefficient-concavity}
  \partial_a^2M(\eps,a,b)
  =\frac{S_0''(a)}{\mathcal A(\eps,b)^2}<0.
\end{equation}

By \eqref{eq:krrs-quotient-boundary}, \(\partial_aL(\eps,a,b,0)=0\); hence, by the
analytic factorization fact above, there is a unique analytic function
\(\mathcal F_1\) satisfying
\begin{equation}\label{eq:krrs-small-block-stationarity}
  \partial_aL(\eps,a,b,\vartheta)
  =\vartheta\mathcal F_1(\eps,a,b,\vartheta).
\end{equation}
Also put
\begin{equation}\label{eq:krrs-cross-block-stationarity}
  \mathcal F_2(\eps,a,b,\vartheta)
  :=\partial_bL(\eps,a,b,\vartheta).
\end{equation}
For fixed \((\eps,\vartheta)\) with \(\vartheta>0\), every interior local
maximizer \((a,b)\) of
\[
  (a,b)\longmapsto\Sigma(\eps,a,b,\vartheta)
\]
satisfies \(\mathcal F_1=\mathcal F_2=0\).  To verify this,
\eqref{eq:krrs-entropy-quotient} gives
\(\partial_a\Sigma=\vartheta^2\mathcal F_1\) and
\(\partial_b\Sigma=\vartheta\mathcal F_2\).

We next show that \((\mathcal F_1,\mathcal F_2)\) vanishes at
\((\eps_0,a_0,z_0,0)\) and that its Jacobian with respect to \((a,b)\) is
invertible there.  These are the hypotheses needed to solve the system
\(\mathcal F_1=\mathcal F_2=0\) locally for \((a,b)\) as analytic functions
of \((\eps,\vartheta)\).  First,
\(\mathcal F_2(\eps_0,a_0,z_0,0)=0\) by
\eqref{eq:krrs-quotient-boundary} and the critical-point property of \(z_0\).  For
fixed \(\eps=\eps_0\) and sufficiently small \(\vartheta>0\), let
\(W_\vartheta\) be the entropy maximizer supplied by
\Cref{thm:krrs-bipodality} for \(\tau=\eps_0^m+\vartheta\).  Denote its bipodal
parameters by \((a,b,q,c)(\vartheta)\), after relabeling so that the first
pode is the one whose size \(c(\vartheta)\) tends to zero as
\(\vartheta\downarrow0\).  By \eqref{eq:krrs-parameter-limits}, they converge to
\((a_0,z_0,\eps_0,0)\), so they eventually lie in the neighborhood
parametrized above.  Since \(W_\vartheta\) has edge density \(\eps_0\),
solving the edge constraint for \(q\) gives
\(q=Q(\eps_0,a,b,c)\).  Its \(H\)-density is
\(\eps_0^m+\vartheta\); hence the local uniqueness assertion in
\eqref{eq:krrs-density-constraint} gives
\(c=C(\eps_0,a,b,\vartheta)\).  Therefore
\(s(W_\vartheta)=\Sigma(\eps_0,a,b,\vartheta)\).  For sufficiently small
\(\vartheta>0\), the three edge-probability parameters lie in \((0,1)\) by
\eqref{eq:krrs-parameter-limits}, while \(c\in(0,1)\) because a null pode would give
zero excess.  Hence, by continuity of \(C\) and \(Q\), every sufficiently
nearby pair \((\tilde a,\tilde b)\), with \(c\) and \(q\) supplied by \(C\)
and \(Q\), yields an admissible bipodal graphon with the same two constraints
and entropy \(\Sigma(\eps_0,\tilde a,\tilde b,\vartheta)\).  Since
\(W_\vartheta\) maximizes entropy among all graphons with these constraints,
\((a,b)(\vartheta)\) is an interior local maximizer of
\(\Sigma(\eps_0,\cdot,\cdot,\vartheta)\) and therefore satisfies
\(\mathcal F_1=\mathcal F_2=0\).  Passage to the limit
\(\vartheta\downarrow0\) gives
\begin{equation}\label{eq:krrs-boundary-stationarity}
  \mathcal F_1(\eps_0,a_0,z_0,0)
  =\mathcal F_2(\eps_0,a_0,z_0,0)=0.
\end{equation}
Equations \eqref{eq:krrs-small-block-stationarity} and \eqref{eq:krrs-quotient-expansion} give
\[
  \mathcal F_1(\eps,a,b,0)=\partial_aM(\eps,a,b),
\]
while \eqref{eq:krrs-cross-block-stationarity} and \eqref{eq:krrs-quotient-boundary} give
\[
  \mathcal F_2(\eps,a,b,0)
  =\frac{\partial_z\psi_d(\eps,b)}{n_d\eps^{m-d}}.
\]
Differentiating these identities at \((\eps_0,a_0,z_0)\), and using
\eqref{eq:krrs-coefficient-concavity} and \eqref{eq:krrs-scalar-nondegeneracy}, gives
\begin{align}
  \partial_a\mathcal F_1(\eps_0,a_0,z_0,0)
  &=\frac{S_0''(a_0)}{\mathcal A(\eps_0,z_0)^2}\ne0,
    \label{eq:krrs-jacobian-aa}\\
  \partial_a\mathcal F_2(\eps_0,a_0,z_0,0)
  &=0,
    \label{eq:krrs-jacobian-ba}\\
  \partial_b\mathcal F_2(\eps_0,a_0,z_0,0)
  &=\frac{\partial_z^2\psi_d(\eps_0,z_0)}
          {n_d\eps_0^{m-d}}\ne0.
    \label{eq:krrs-jacobian-bb}
\end{align}
The Jacobian of \((\mathcal F_1,\mathcal F_2)\) with respect to \((a,b)\)
is therefore triangular with nonzero diagonal.  The analytic implicit
function theorem now gives an open interval \(U_{\mathrm{IFT}}\ni\eps_0\),
a number \(\Delta_{\mathrm{IFT}}>0\), an open neighborhood
\(V_{\mathrm{IFT}}\) of \((a_0,z_0)\), and real-analytic functions
\begin{equation}\label{eq:krrs-stationary-densities}
  a_*(\eps,\vartheta),\qquad b_*(\eps,\vartheta)
  \quad
  ((\eps,\vartheta)\in
    U_{\mathrm{IFT}}\times(-\Delta_{\mathrm{IFT}},
                             \Delta_{\mathrm{IFT}}))
\end{equation}
taking values in \(V_{\mathrm{IFT}}\), with
\[
  (a_*(\eps_0,0),b_*(\eps_0,0))=(a_0,z_0).
\]
For each \((\eps,\vartheta)\) in this rectangle,
\[
  \mathcal F_i\bigl(\eps,a_*(\eps,\vartheta),
    b_*(\eps,\vartheta),\vartheta\bigr)=0
  \qquad(i=1,2),
\]
and \((a_*(\eps,\vartheta),b_*(\eps,\vartheta))\) is the unique solution in
\(V_{\mathrm{IFT}}\) of this system.  Define
\begin{align}
  c_*(\eps,\vartheta)
  &:=C\bigl(\eps,a_*(\eps,\vartheta),
                 b_*(\eps,\vartheta),\vartheta\bigr),\label{eq:krrs-block-size-map}\\
  q_*(\eps,\vartheta)
  &:=Q\bigl(\eps,a_*(\eps,\vartheta),
                 b_*(\eps,\vartheta),c_*(\eps,\vartheta)\bigr).
                 \label{eq:krrs-large-block-map}
\end{align}
After shrinking \(U_{\mathrm{IFT}}\) and \(\Delta_{\mathrm{IFT}}\) so that
these compositions are defined, all four functions are real-analytic on the
two-sided rectangle
\(U_{\mathrm{IFT}}\times(-\Delta_{\mathrm{IFT}},\Delta_{\mathrm{IFT}})\).
At \(\vartheta=0\), equations \eqref{eq:krrs-density-constraint} and
\eqref{eq:krrs-edge-constraint} give
\[
  c_*(\eps,0)=0,
  \qquad
  q_*(\eps,0)=\eps.
\]
Moreover, \(\mathcal F_2(\eps,a_*(\eps,0),b_*(\eps,0),0)=0\) and
\eqref{eq:krrs-quotient-boundary} show that
\(b_*(\eps,0)\) is a critical point of
\(z\mapsto\psi_d(\eps,z)\).  The uniqueness assertion in
\Cref{thm:krrs-cross-density} shows that this critical point is
unique, so, after shrinking \(U_{\mathrm{IFT}}\),
\begin{equation}\label{eq:krrs-cross-block-limit}
  b_*(\eps,0)=\zeta_d(\eps)
  \qquad(\eps\in U_{\mathrm{IFT}}).
\end{equation}
In particular, this also proves directly that \(\zeta_d\) is real-analytic
locally away from \(\eps_*\).

\paragraph{4. Identification with the entropy-maximizer parameters.}
Let \(\Omega^+\) denote the domain on which
\Cref{thm:krrs-bipodality} supplies jointly real-analytic bipodal optimizer
parameters, expressed in the coordinates \((\eps,\vartheta)\) with
\(\tau=\eps^m+\vartheta\).  For each nonexceptional \(\eps\), this domain is
specified by
\(\eps^m<\tau<\tau_0(\eps)\) for a threshold \(\tau_0(\eps)>\eps^m\)
depending on \(\eps\), so the fiber of
\(\Omega^+\) over each nonexceptional \(\eps\) is an interval
\begin{equation}\label{eq:krrs-domain-fibres}
  \{\vartheta:(\eps,\vartheta)\in\Omega^+\}
  =(0,h(\eps)),
  \qquad
  h(\eps):=\tau_0(\eps)-\eps^m>0.
\end{equation}
Since the theorem states that the parameter map is
real-analytic on \(\Omega^+\), this domain is open (real-analytic maps are,
by definition, defined on open sets).

Choose \(\vartheta^\sharp\in(0,h(\eps_0))\).  Openness at
\((\eps_0,\vartheta^\sharp)\) gives an open interval
\(U_{\mathrm{KRR}}\ni\eps_0\) such that
\((\eps,\vartheta^\sharp)\in\Omega^+\) for every
\(\eps\in U_{\mathrm{KRR}}\).  The interval property
\eqref{eq:krrs-domain-fibres} then implies
\begin{equation}\label{eq:krrs-uniform-domain}
  U_{\mathrm{KRR}}\times(0,\vartheta^\sharp)
  \subseteq\Omega^+.
\end{equation}
Thus, the same interval \(0<\vartheta<\vartheta^\sharp\) lies in the domain of
\Cref{thm:krrs-bipodality} for every \(\eps\in U_{\mathrm{KRR}}\).  This
conclusion uses only the theorem statement and not any localization argument
from its proof.

Shrink to an open interval \(U\ni\eps_0\) whose closure is contained in
\[
  U_{\mathrm{IFT}}\cap U_{\mathrm{KRR}}
  \cap\bigl((0,1)\setminus\{\eps_*\}\bigr),
\]
and choose
\[
  0<\Delta<\min\{\Delta_{\mathrm{IFT}},\vartheta^\sharp\}.
\]
By \Cref{thm:krrs-bipodality}, throughout
\(U\times(0,\Delta)\) the entropy maximizer is unique up to relabeling and
bipodal.  For each \((\eps,\vartheta)\) in this set, let
\(W^+_{\eps,\vartheta}\) be a representative of this maximizer, with its
podes labeled so that the first pode has size \(c^+\) and
\(c^+\to0\) as \(\vartheta\downarrow0\).  The corresponding parameter map
\begin{equation}\label{eq:krrs-optimizer-map}
  (a^+,b^+,q^+,c^+)(\eps,\vartheta)
\end{equation}
is real-analytic.

It remains to identify the optimizer parameter map
\eqref{eq:krrs-optimizer-map} with the two-sided analytic solution of the
stationarity equations constructed in
\eqref{eq:krrs-stationary-densities}--\eqref{eq:krrs-large-block-map}.  At \(\eps=\eps_0\), the
one-sided limits in \Cref{thm:krrs-bipodality} and
\eqref{eq:krrs-parameter-limits} show that the parameters in
\eqref{eq:krrs-optimizer-map} lie in the local solution neighborhoods
for all sufficiently small \(\vartheta>0\).  Choose one such
\(\vartheta_1\in(0,\Delta)\).  The three edge-probability parameters have
limits \(a_0,z_0,\eps_0\in(0,1)\), and \(c^+(\eps_0,\vartheta_1)\in(0,1)\)
because a bipodal graphon with a null pode is constant and has excess
\(0\); after taking \(\vartheta_1\) sufficiently small, all four parameters
therefore lie in \((0,1)\).  By continuity, the same is true on a nonempty
open neighborhood \(\mathcal O\subset U\times(0,\Delta)\) of
\((\eps_0,\vartheta_1)\), and all parameters there remain in those
neighborhoods.

For every \((\eps,\vartheta)\in\mathcal O\), we repeat the constraint and
local-maximality argument used to prove \eqref{eq:krrs-boundary-stationarity}.  The edge
constraint gives \(q^+=Q(\eps,a^+,b^+,c^+)\), and the local uniqueness in
\eqref{eq:krrs-density-constraint} gives
\(c^+=C(\eps,a^+,b^+,\vartheta)\).  Since
\(W^+_{\eps,\vartheta}\) maximizes entropy, it is locally maximal among the
nearby bipodal competitors supplied by \(C\) and \(Q\).  Hence
\(\mathcal F_1=\mathcal F_2=0\) at
\((\eps,a^+,b^+,\vartheta)\), and the local uniqueness obtained above gives
\begin{equation}\label{eq:krrs-map-agreement}
  (a^+,b^+,q^+,c^+)
  =(a_*,b_*,q_*,c_*)
  \qquad\text{on }\mathcal O.
\end{equation}
Both sides of \eqref{eq:krrs-map-agreement} are real-analytic on
\(U\times(0,\Delta)\), which is a product of intervals and hence connected.
The one-sided limits in \Cref{thm:krrs-bipodality} hold as
\(\vartheta\downarrow0\) for each fixed \(\eps\), but are not asserted to be
uniform in \(\eps\).  Thus, the argument based on these limits initially
identifies the two parameter maps only on \(\mathcal O\).  The identity
theorem for real-analytic functions extends this equality to all of the
connected set \(U\times(0,\Delta)\).  Consequently,
\((a_*,b_*,q_*,c_*)\) is the desired two-sided real-analytic extension of the
optimizer parameter map \eqref{eq:krrs-optimizer-map}.

Since \(\overline U\) avoids \(\eps_*\), the fixed-point assertion in
\Cref{thm:krrs-cross-density} gives
\(\zeta_d(\eps)\ne\eps\), and hence
\(\mathcal D_\eps(\zeta_d(\eps))>0\), for every \(\eps\in\overline U\).
Finally, since \(C(\eps,a,b,0)\equiv0\), the partial derivatives
\(\partial_aC\) and \(\partial_bC\) vanish at \(\vartheta=0\).
Differentiating \eqref{eq:krrs-block-size-map} with respect to \(\vartheta\) at
\(\vartheta=0\) therefore yields
\[
  \partial_\vartheta c_*(\eps,0)
  =\partial_\vartheta C\bigl(
     \eps,a_*(\eps,0),b_*(\eps,0),0\bigr).
\]
Equations \eqref{eq:krrs-block-size-expansion} and \eqref{eq:krrs-cross-block-limit} now give
\begin{equation}\label{eq:krrs-block-size-derivative}
  \partial_\vartheta c_*(\eps,0)
  =\frac{1}{n_d\eps^{m-d}
       \mathcal D_\eps(\zeta_d(\eps))}>0.
\end{equation}
Moreover, \(c_*\) is analytic on an open rectangle containing the compact
set \(\overline U\times[-\Delta,\Delta]\) and vanishes on
\(\{\vartheta=0\}\), so \(\partial_\vartheta c_*\) is bounded there, and
the mean value theorem gives, uniformly for \(\eps\in U\),
\[
  c_*(\eps,\vartheta)=O_{H,U}(\vartheta)
  \qquad(\vartheta\to0).
\]
Finally, \(c_*(\eps,\vartheta)\) takes values in \((0,1)\) when
\(\eps\in U\) and \(0<\vartheta<\Delta\): this follows, after shrinking
\(\Delta\) once more, from \eqref{eq:krrs-block-size-derivative} and
\(c_*(\eps,0)=0\).

Together with \eqref{eq:krrs-cross-block-limit} and the boundary identities for
\(q_*\) and \(c_*\), this proves every assertion of
\Cref{thm:krrs-analytic-extension}.
\end{proof}
